\chapter{Overview}
\label{chap:overview}
% archon:covers MR4583777TatesThesisInTheDeRhamSetting.lean MR4583777TatesThesisInTheDeRhamSetting/Basic.lean

This chapter provides the complete informal blueprint for the formalization of
Hilburn--Raskin \cite{hilburn-raskin} (arXiv:2107.11325, MR4583777).
The main theorem is a canonical equivalence of DG categories
\[
  \Delta : D^!(\LLA) \simeq \IndCohStar(\cY)
\]
where $\LLA$ is the algebraic loop space of $\BA^1$ and $\cY$ is a moduli
space of rank 1 de Rham local systems on the punctured formal disc with a
flat section. We follow the paper's structure: \S\ref{sec:loop-spaces} covers
algebraic loop spaces and local systems; \S\ref{sec:Y} covers the space $\cY$;
\S\ref{sec:indcoh} covers ind-coherent sheaves; \S\ref{sec:dmod} covers
D-modules and the Weyl algebra; \S\ref{sec:main-results} states the main theorems.

\section{Geometry of the spectral side}
\label{sec:loop-spaces}
\label{sec:Y}

\begin{definition}[Loop space]
  \label{def:loop-space}
  \lean{MR4583777TatesThesisInTheDeRhamSetting.TODO.loopSpace}
  \uses{}
  % SOURCE: Hilburn-Raskin §2.1.1 (loop functor)
  % SOURCE QUOTE: "For a prestack $Z$, $\fL Z \in \PreStk \coloneqq \TwoHom(\AffSch^{op},\Gpd)$ denote the functor: $\Spec(A) \mapsto \Hom_{\AffSch}(\Spec(A((t))),Z)$."
  For a prestack $Z \in \PreStk$, the \emph{loop space} $\mathfrak{L}Z$
  is the prestack defined on $\AffSch$ by
  \[
    \Spec(A) \;\mapsto\; \Hom_{\AffSch}(\Spec(A((t))),\, Z).
  \]
  For $Z = \BA^1$ this gives $\LLA(\Spec A) = A((t))$, so $\LLA$ is the
  ind-scheme whose $A$-points are formal Laurent series with $A$-coefficients.

  \textit{Source: \cite{hilburn-raskin}, \S2.1.1.}
\end{definition}

\begin{definition}[Positive loop space]
  \label{def:positive-loop-space}
  \lean{MR4583777TatesThesisInTheDeRhamSetting.TODO.posLoopSpace}
  \uses{}
  % SOURCE: Hilburn-Raskin §2.1.1 (positive loop functor)
  % SOURCE QUOTE: "Similarly, we let $\fL^+ Z$ denote the functor: $\Spec(A) \mapsto \Hom_{\AffSch}(\Spec(A[[t]]),Z)$."
  The \emph{positive loop space} $\mathfrak{L}^+ Z$ for $Z \in \PreStk$
  is defined by
  \[
    \Spec(A) \;\mapsto\; \Hom_{\AffSch}(\Spec(A[[t]]),\, Z).
  \]
  For $Z = \BA^1$ this gives $\LLAp(\Spec A) = A[[t]]$, so $\LLAp$ is the
  pro-scheme whose $A$-points are formal power series with $A$-coefficients.

  \textit{Source: \cite{hilburn-raskin}, \S2.1.1.}
\end{definition}

\begin{definition}[Truncated positive loop space]
  \label{def:truncated-positive-loop-space}
  \lean{MR4583777TatesThesisInTheDeRhamSetting.TODO.truncPosLoopSpace}
  \uses{def:positive-loop-space}
  % SOURCE: Hilburn-Raskin §2.1.1 (truncated positive loop)
  % SOURCE QUOTE: "For $n \geq 0$, we let $\fL_n^+$ denote the functor: $\Spec(A) \mapsto \Hom_{\AffSch}(\Spec(A[[t]]/t^n),Z)$."
  For $n \geq 0$ and $Z \in \PreStk$, the \emph{truncated positive loop space}
  $\LLAn{n} := \mathfrak{L}_n^+ Z$ is defined by
  \[
    \Spec(A) \;\mapsto\; \Hom_{\AffSch}(\Spec(A[[t]]/t^n),\, Z).
  \]
  For $Z = \BA^1$ this gives $\LLAn{n} \cong \BA^n$ via the identification
  $A[[t]]/t^n \cong A^n$ (coefficients $a_0,\ldots,a_{n-1}$).

  \textit{Source: \cite{hilburn-raskin}, \S2.1.1.}
\end{definition}

\begin{definition}[Local systems on the punctured formal disc]
  \label{def:locsys-gm}
  \lean{MR4583777TatesThesisInTheDeRhamSetting.TODO.locSysGm}
  \uses{def:loop-space}
  % SOURCE: Hilburn-Raskin §\ref{ss:rk1-ls}
  % SOURCE QUOTE: "We have a map $\fL \bG_m \to \fL \bA^1 dt$ sending $f \in \fL \bG_m$ to $d\log(f)$... We define $\LocSys_{\bG_m}$ as the stack quotient $\fL \bA^1 dt/\fL \bG_m$." / "for a point $\omega \in \fL \bA^1 dt$, we consider $\nabla = d - \omega$ as the corresponding rank $1$ connection."
  The moduli stack of \emph{rank 1 de Rham local systems on the punctured
  formal disc} is the stack quotient
  \[
    \LocSysGm \;:=\; \LLAdt \,/\, \LLGm
  \]
  where $\LLAdt$ denotes the loop space of $1$-forms $\BA^1 dt$, and $\LLGm$
  acts \emph{additively} on $\LLAdt$ through the homomorphism
  $d\log : \LLGm \to \LLAdt$, i.e.\ by gauge transformations
  $g \cdot \omega = \omega + d\log(g)$. A point $\omega$ corresponds to the
  connection $\nabla = d - \omega$ on the trivial line bundle.

  \textit{Source: \cite{hilburn-raskin}, \S\texttt{ss:rk1-ls}.}
\end{definition}

\begin{definition}[Polar forms]
  \label{def:polar-forms}
  \lean{MR4583777TatesThesisInTheDeRhamSetting.TODO.polarForms}
  \uses{def:loop-space, def:positive-loop-space}
  % SOURCE: Hilburn-Raskin §\ref{ss:rk1-ls}
  % SOURCE QUOTE: "Define $\fL^{pol} \bA^1 dt$ as the quotient $\fL \bA^1 dt / \fL^+\bA^1 dt$, the quotient being with respect to the additive structure. In other words, $\fL^{pol} \bA^1 dt$ parametrizes polar parts of differential forms on the disc. Note that $\fL^{pol} \bA^1 dt$ is an indscheme of ind-finite type."
  The space of \emph{polar forms} is the \emph{additive} quotient
  \[
    \fL^{pol}\BA^1 dt \;:=\; \LLAdt / \LLAp dt,
  \]
  an indscheme of ind-finite type. The \emph{polar part} of
  $f(t)dt \in k((t))dt$ is its class in
  $k((t))dt / k[[t]]dt \cong t^{-1}k[t^{-1}]\,dt$.

  \textit{Source: \cite{hilburn-raskin}, \S\texttt{ss:rk1-ls}.}
\end{definition}

\begin{proposition}[Logarithmic decomposition of $\LocSysGm$]
  \label{prop:locsys-log-decomp}
  \lean{MR4583777TatesThesisInTheDeRhamSetting.TODO.locSysLogDecomp}
  \uses{def:locsys-gm, def:polar-forms}
  % SOURCE: Hilburn-Raskin Proposition \ref{p:locsys-log}
  % SOURCE QUOTE: "The map: $\LocSys_{\bG_m,\log} \to \fL^{pol} \bA^1 dt \times \bB \bG_m$ is an isomorphism."
  Let $\LocSysGm_{\log} := \LLAdt/\LLGmp$ be the logarithmic local systems
  (line bundles on the disc with a connection on the punctured disc), with
  $\LocSysGm = \LocSysGm_{\log}/\GrGm$. The polar-part map
  $\on{Pol} : \LocSysGm_{\log} \to \fL^{pol}\BA^1 dt$ together with the
  fiber-at-origin map $\LocSysGm_{\log}\to \bB\Gm$ give an isomorphism
  \[
    \LocSysGm_{\log} \;\xrightarrow{\;\sim\;}\; \fL^{pol}\BA^1 dt \times \bB\Gm.
  \]

  \textit{Source: \cite{hilburn-raskin}, Proposition \texttt{p:locsys-log}.}
\end{proposition}

% SOURCE QUOTE PROOF: (read from references/MR4583777-de-rham-tate.tex)
% "The map $\fL^+ \bG_m \xar{d\log,\ev} \fL^+ \bA^1 dt \times \bG_m$
% is easily seen to be an isomorphism. The result is immediate from here."
\begin{proof}
  \uses{def:locsys-gm, def:polar-forms}
  The map $d\log : \LLGmp \to \LLAp dt$ and the evaluation-at-$0$ map
  $\on{ev} : \LLGmp \to \Gm$ combine to give an isomorphism
  $(\LLGmp \xrightarrow{(d\log,\,\on{ev})} \LLAp dt \times \Gm)$,
  which is verified on $A$-points: $(g = \sum a_i t^i) \mapsto
  (g^{-1}dg,\, a_0)$ is invertible by the formal power series inverse.
  The logarithmic local systems $\LocSysGm_{\log}$ are then identified
  with $\mathfrak{L}^{\mathrm{pol}}\BA^1 dt \times B\Gm$ via this
  isomorphism and the quotient construction.
\end{proof}

\begin{proposition}[Decomposition of $\LocSysGm$]
  \label{prop:locsys-decomp}
  \lean{MR4583777TatesThesisInTheDeRhamSetting.TODO.locSysDecomp}
  \uses{prop:locsys-log-decomp}
  % SOURCE: Hilburn-Raskin Proposition \ref{p:locsys}
  % SOURCE QUOTE: "There is an isomorphism $\LocSys_{\bG_m} \simeq \bA^1/\bZ \times \Ker(\Res)_{dR} \times \bB \bG_m$ fitting into a commutative diagram [with Prop \ref{p:locsys-log}]."
  There is an isomorphism of prestacks
  \[
    \LocSysGm \;\simeq\; \BA^1/\BZ \;\times\; \on{Ker}(\Res)_{dR}
    \;\times\; \bB\Gm.
  \]
  Here $\BA^1/\BZ$ classifies the residue (a $k$-point plus a monodromy),
  $\on{Ker}(\Res)_{dR}$ is the de Rham local system with trivial residue,
  and $\bB\Gm$ is the classifying stack of $\Gm$.

  \textit{Source: \cite{hilburn-raskin}, Proposition \texttt{p:locsys}.}
\end{proposition}

% SOURCE QUOTE PROOF: (read from references/MR4583777-de-rham-tate.tex)
% "There is a canonical homomorphism $\Gr_{\bG_m} \to \bZ$ given by minus
% the valuation map. This map splits canonically as well:
% $\bZ = \Gr_{\bG_m}^{red}$.
% Therefore, we obtain a canonical product decomposition:
% $\Gr_{\bG_m}^{red} = \bZ \times \Gr_{\bG_m}^{\circ}$
% for $\Gr_{\bG_m}^{\circ}$ the connected component of the identity.
% [The map $d\log:\Gr_{\bG_m}^{\circ} \to \Ker(\Res)$] is easily seen to
% induce an isomorphism between $\Gr_{\bG_m}^{\circ}$ and the formal group
% $\Ker(\Res)_0^{\wedge}$ of $\Ker(\Res)$. This gives the claim."
\begin{proof}
  \uses{prop:locsys-log-decomp}
  From Proposition~\ref{prop:locsys-log-decomp}, the logarithmic local
  systems $\LocSysGm_{\log} \cong \mathfrak{L}^{\mathrm{pol}}\BA^1 dt
  \times B\Gm$ with residue decomposition
  $\mathfrak{L}^{\mathrm{pol}}\BA^1 dt = \BA^1 \times \on{Ker}(\Res)$.
  The Grassmannian $\on{Gr}_{\Gm}$ decomposes as $\BZ \times \on{Gr}_{\Gm}^{\circ}$
  via the (minus-) valuation map, where the $\BZ$-factor records the
  residue (contributing the $\BA^1/\BZ$ factor via the exponential) and
  $d\log$ identifies $\on{Gr}_{\Gm}^{\circ}$ with the formal group
  $\on{Ker}(\Res)_0^{\wedge}$. Quotienting by the $\LLGm$-action gives
  $\LocSysGm \cong \BA^1/\BZ \times \on{Ker}(\Res)_{dR} \times B\Gm$.
\end{proof}

\begin{definition}[The prestack $\cY'$]
  \label{def:Y-prime}
  \lean{MR4583777TatesThesisInTheDeRhamSetting.TODO.Yprime}
  \uses{def:loop-space}
  % SOURCE: Hilburn-Raskin §\ref{ss:y-defin}
  % SOURCE QUOTE: "$\sY^{\prime} \coloneqq \on{Eq}\big(\fL \bA^1 \times \fL \bA^1 dt \underset{(g,\omega) \mapsto g\omega}{\overset{d\circ \pi_1}{\rightrightarrows}} \fL \bA^1 dt\big) \in \PreStk$" / "$\sY^{\prime}$ by definition parametrizes pairs $(g \in \fL\bA^1, \omega \in \fL \bA^1 dt)$ with $dg = g\omega$, which we can rewrite as $\nabla g = 0$ for $\nabla \coloneqq d - \omega$."
  Define the prestack
  \[
    \cY' \;:=\; \Eq\!\Big(\LLA \times \LLAdt
    \;\underset{(g,\omega)\mapsto g\omega}{\overset{d \circ \pi_1}{\rightrightarrows}}\;
    \LLAdt\Big) \;\in\; \PreStk.
  \]
  Concretely, the $A$-points of $\cY'$ are pairs
  $(g, \omega) \in A((t)) \times A((t))dt$ satisfying $dg = g\omega$, i.e.\
  $\nabla g = 0$ for the connection $\nabla := d - \omega$; so $\cY'$ encodes
  a connection on the trivial bundle together with a flat section.
  This is a \emph{derived} equalizer (a priori a DG indscheme).

  \textit{Source: \cite{hilburn-raskin}, \S\texttt{ss:y-defin}.}
\end{definition}

\begin{definition}[The space $\cY$]
  \label{def:Y}
  \lean{MR4583777TatesThesisInTheDeRhamSetting.TODO.Y}
  \uses{def:Y-prime, def:locsys-gm}
  % SOURCE: Hilburn-Raskin §\ref{ss:y-defin}
  % SOURCE QUOTE: "$\sY \coloneqq \sY^{\prime}/\fL \bG_m.$" / action: "$(f \in \fL \bG_m,(g,\omega) \in \sY^{\prime}) \mapsto (fg,\omega+d\log(f))$."
  The \emph{de Rham moduli space} is the stack quotient
  \[
    \cY \;:=\; \cY' / \LLGm
  \]
  where $\LLGm$ acts on $\cY'$ by gauge transformations
  $f \cdot (g,\omega) = (fg,\; \omega + d\log(f))$ (the action is made
  rigorous in \S\,\texttt{ss:action-rigorous} of \cite{hilburn-raskin}).
  There is a canonical map $\pi : \cY \to \LocSysGm$ recording $\omega$.

  \textit{Source: \cite{hilburn-raskin}, \S\texttt{ss:y-defin}.}
\end{definition}

\begin{theorem}[$\cY$ is classical]
  \label{thm:Y-classical}
  \lean{MR4583777TatesThesisInTheDeRhamSetting.TODO.YIsClassical}
  \uses{def:Y, lem:flat-multiplication}
  % SOURCE: Hilburn-Raskin Theorem \ref{t:cl}
  % SOURCE QUOTE: "$\sY$ is a classical prestack." / "We formally deduce the same for $\sY^{\prime}$."
  $\cY$ is a classical prestack. In particular, the \emph{derived} equalizer
  defining $\cY'$ is classical (one formally deduces classicality of $\cY'$
  from that of $\cY$).

  \textit{Source: \cite{hilburn-raskin}, Theorem \texttt{t:cl}.}
\end{theorem}

% SOURCE QUOTE PROOF: (read from references/MR4583777-de-rham-tate.tex)
% "By Corollary \ref{c:yn-log-cl},
% $\sY^{\leq n} = \sY_{\log}^{\leq n}/\Gr_{\bG_m}^{\leq n}$
% is a classical (non-algebraic) prestack.
% We deduce the same for $\sY = \colim_n \sY^{\leq n}$."
\begin{proof}
  \uses{def:Y, def:Z-leq-n, lem:flat-multiplication}
  By Corollary~c:yn-log-cl in \cite{hilburn-raskin}, each truncated space
  $\cY^{\leq n} = \cY_{\log}^{\leq n}/\on{Gr}_{\Gm}^{\leq n}$ is a classical
  (non-algebraic) prestack. This rests on the flatness results of
  \S2.6 (Lemma~\ref{lem:flat-multiplication}), which show the truncated
  spaces $\cZ^{\leq n}$ are classical. Since $\cY = \varinjlim_n \cY^{\leq n}$
  is a filtered colimit of classical prestacks, $\cY$ is classical as well.
\end{proof}

\begin{definition}[Truncated spaces $\cZ^{\leq n}$ and $\cY^{\leq n}$]
  \label{def:Z-leq-n}
  \lean{MR4583777TatesThesisInTheDeRhamSetting.TODO.ZleqN}
  \uses{def:Y, def:positive-loop-space, def:locsys-gm}
  % SOURCE: Hilburn-Raskin §\ref{ss:y-trun} (Intermediate spaces; truncated $\sZ^{\leq n}$, $\sY^{\leq n}$)
  % SOURCE QUOTE: "$\sZ^{\leq n} \coloneqq \on{Eq}\big(\fL^+ \bA^1 \times \fL^{\leq n} \bA^1 dt \underset{(g,\omega) \mapsto g\omega}{\overset{d\circ \pi_1}{\rightrightarrows}} \fL^{\leq n} \bA^1 dt\big)/\fL^+ \bG_m$"
  For $n \geq 0$, the \emph{$n$-th truncated moduli space} is the prestack
  \[
    \cZ^{\leq n} \;:=\; \on{Eq}\!\Big(\LLAp \times \LLApol{n}dt
    \;\underset{(g,\omega)\mapsto g\omega}{\overset{d \circ \pi_1}{\rightrightarrows}}\;
    \LLApol{n}dt\Big)\Big/ \LLGmp.
  \]
  Here the \emph{first} factor is the \emph{positive} loop space $\LLAp$
  (power series $A[[t]]$), the forms $\LLApol{n}dt := \fL^{\leq n}\BA^1 dt$
  range over $1$-forms with \emph{poles of order at most $n$}, and the
  quotient is by the \emph{positive} loop group $\LLGmp$ (not by all of
  $\LLGm$). The $\LLGmp$-action is constructed as in
  \S\,\texttt{ss:action-rigorous} of \cite{hilburn-raskin}.
  There are closed immersions $\cZ^{\leq n} \xrightarrow{\iota} \cZ^{\leq n}$
  (the shift map $\iota$), whose colimit is the logarithmic truncated space
  $\cY_{\log}^{\leq n} = \varinjlim_\iota \cZ^{\leq n}$, and one sets
  $\cY^{\leq n} := \cY_{\log}^{\leq n}/\on{Gr}_{\Gm}^{\leq n}$.
  The full space is the colimit $\cY = \varinjlim_n \cY^{\leq n}$.

  \textit{Source: \cite{hilburn-raskin}, \S\texttt{ss:y-trun}.}
\end{definition}

\begin{proposition}[$\cZ^{\leq n}$ is an algebraic stack]
  \label{prop:Z-leq-n-algebraic}
  \lean{MR4583777TatesThesisInTheDeRhamSetting.TODO.ZleqNAlgebraic}
  \uses{def:Z-leq-n, lem:flat-multiplication}
  % SOURCE: Hilburn-Raskin Proposition \ref{p:zn-algebraic}
  % SOURCE QUOTE: "The prestack $\sZ^{\leq n}$ is a (DG) algebraic stack."
  For each $n \geq 0$, the prestack $\cZ^{\leq n}$ is a (DG) algebraic stack.

  \textit{Source: \cite{hilburn-raskin}, Proposition \texttt{p:zn-algebraic}.}
\end{proposition}

% SOURCE QUOTE PROOF: (read from references/MR4583777-de-rham-tate.tex)
% "By definition, we have:
% $\on{Eq}\big(\fL^+ \bA^1 \times \fL^{\leq n} \bA^1 dt
% \underset{(g,\omega) \mapsto g\omega}{\overset{d\circ \pi_1}{\rightrightarrows}}
% \fL^{\leq n} \bA^1 dt\big) =
% \fL^{\leq n} \bA^1 dt \underset{\LocSys_{\bG_m,\log}^{\leq n}}{\times}
% \sZ^{\leq n}$.
% Therefore, $\sZ^{\leq n} \to \LocSys_{\bG_m,\log}^{\leq n}$ is
% an affine morphism. As $\LocSys_{\bG_m,\log}^{\leq n}$ is an algebraic stack
% by Proposition \ref{p:locsys-trun}, we obtain the result."
\begin{proof}
  \uses{def:Z-leq-n, lem:flat-multiplication, prop:locsys-log-decomp}
  From the definition of $\cZ^{\leq n}$ as a fiber product, the equalizer
  formula shows that the map $\cZ^{\leq n} \to \LocSysGm_{\log}^{\leq n}$
  is affine. The truncated logarithmic local systems
  $\LocSysGm_{\log}^{\leq n}$ form an algebraic stack
  (Proposition~p:locsys-trun in \cite{hilburn-raskin}, which follows from
  Proposition~\ref{prop:locsys-log-decomp} and the algebraicity of
  $\LLAn{n}$). A morphism that is affine over an algebraic stack remains
  algebraic, so $\cZ^{\leq n}$ is an algebraic stack. The flatness results
  of \S2.6 (Lemma~\ref{lem:flat-multiplication}) are what make the
  fiber-product/equalizer presentation behave well.
\end{proof}

\begin{definition}[Truncated convolution map $\mu_n$]
  \label{def:trunc-conv}
  \lean{MR4583777TatesThesisInTheDeRhamSetting.TODO.truncConvMap}
  \uses{def:truncated-positive-loop-space}
  % SOURCE: Hilburn-Raskin §\ref{ss:flatness}, Lemma \ref{l:flat-mult}
  % SOURCE QUOTE: "The map: $\mu_n:\bA^n \times \bA^n = \bA^{2n} \to \bA^n$ $(a_0,\ldots,a_{n-1},b_0,\ldots,b_{n-1}) \mapsto (a_0b_0,a_0b_1+a_1b_0,\ldots,\sum_{i = 0}^{n-1} a_ib_{n-1-i})$"
  The \emph{truncated convolution map} $\mu_n : \BA^n \times \BA^n \to \BA^n$
  corresponds to multiplication of truncated power series:
  \[
    \mu_n\bigl((a_0,\ldots,a_{n-1}),\,(b_0,\ldots,b_{n-1})\bigr)
    \;:=\; \Big(\sum_{i+j=0} a_i b_j,\; \sum_{i+j=1} a_i b_j,\; \ldots,\;
    \sum_{i+j=n-1} a_i b_j\Big).
  \]
  This is the scheme map $\LLAn{n} \times_k \LLAn{n} \to \LLAn{n}$ induced by
  multiplication $k[[t]]/t^n \otimes_k k[[t]]/t^n \to k[[t]]/t^n$.

  \textit{Source: \cite{hilburn-raskin}, Lemma \texttt{l:flat-mult}.}
\end{definition}

\begin{lemma}[Flatness of truncated convolution]
  \label{lem:flat-multiplication}
  \lean{MR4583777TatesThesisInTheDeRhamSetting.TODO.flatMultiplication}
  \uses{def:trunc-conv}
  % SOURCE: Hilburn-Raskin Lemma \ref{l:flat-mult}
  % SOURCE QUOTE: "The map: $\mu_n:\bA^n \times \bA^n = \bA^{2n} \to \bA^n$ $(a_0,\ldots,a_{n-1},b_0,\ldots,b_{n-1}) \mapsto (a_0b_0,a_0b_1+a_1b_0,\ldots,\sum_{i = 0}^{n-1} a_ib_{n-1-i})$ is flat."
  The truncated convolution map $\mu_n : \BA^{2n} \to \BA^n$ is flat.

  \textit{Source: \cite{hilburn-raskin}, Lemma \texttt{l:flat-mult}.}
\end{lemma}

% SOURCE QUOTE PROOF: (read from references/MR4583777-de-rham-tate.tex)
% "This map is defined by a family of homogeneous polynomials (of degree $2$).
% Therefore, it suffices to show that the fiber $Z \coloneqq \mu_n^{-1}(0)$
% over $0$ is equidimensional with the expected dimension $2n - n = n$.
% First, for $0 \leq m \leq n$, let $Z_m \subset Z$ be the locally closed
% subscheme where $a_0 = a_1 = \ldots = a_{m-1} = 0$ and $a_m \neq 0$,
% the last condition being considered as vacuous for $m = n$.
% We claim that $\dim(Z_m) = n$ for all $m$."
\begin{proof}
  \uses{def:trunc-conv}
  The map $\mu_n$ is defined by degree-$2$ homogeneous polynomials, so
  by a standard homogeneity argument it suffices to show that the central
  fiber $Z := \mu_n^{-1}(0) \subset \BA^{2n}$ is equidimensional of
  dimension $n$. Stratify $Z$ by locally closed subschemes
  $Z_m$ ($0 \leq m \leq n$) where $a_0 = \cdots = a_{m-1} = 0$ and
  $a_m \neq 0$ (vacuous for $m = n$). For $m < n$, the equations
  $\sum_{i=0}^j a_{m+i} b_{j-i} = 0$ ($j = 0,\ldots,n-m-1$) allow one
  to express $b_0,\ldots,b_{n-m-1}$ uniquely in terms of $a_m,\ldots,a_{n-1}$
  and $b_{n-m},\ldots,b_{n-1}$, giving an isomorphism
  $Z_m \cong (\BA^1 \setminus 0) \times \BA^{n-m-1} \times \BA^m$,
  which has dimension $n$. Hence $Z$ is equidimensional of dimension $n$,
  so $\mu_n$ is flat.
\end{proof}

\begin{definition}[The coordinate ring $A_n$ of $\widetilde{\cZ}^{\leq n}$]
  \label{def:An-coords}
  \lean{MR4583777TatesThesisInTheDeRhamSetting.TODO.AnCoordinateRing}
  \uses{def:Z-leq-n}
  % SOURCE: Hilburn-Raskin §\ref{ss:coords} (Definition of $A_n$) and Prop \ref{p:zn-coords}
  % SOURCE QUOTE: "Define a classical commutative ring $A_n$ by taking (infinitely many) generators $a_0,a_1,a_2,\ldots$ and $b_{-1},\ldots,b_{-n}$ and (infinitely) relations coming from equating Taylor series coefficients of the formal Laurent series: $\sum_{i=0}^{\infty} ia_it^{i-1} = \sum_{i=0}^{\infty} a_it^i \cdot \sum_{i=-n}^{-1} b_{-i}t^{-i}$."
  Define a classical commutative ring $A_n$ by taking the
  \emph{infinitely many} generators $a_0, a_1, a_2, \ldots$ together with
  $b_{-1}, \ldots, b_{-n}$, modulo the \emph{infinitely many} relations
  obtained by equating Taylor coefficients of the formal Laurent series
  \[
    \sum_{i=0}^{\infty} i\,a_i t^{i-1}
    \;=\;
    \Big(\sum_{i=0}^{\infty} a_i t^i\Big)\cdot
    \Big(\sum_{i=-n}^{-1} b_{-i} t^{-i}\Big).
  \]
  Here $a_i$ encode a flat section $s = \sum_{i\geq 0} a_i t^i$ and the
  $b_{-i}$ encode the polar connection $\nabla = d - \sum_{i=1}^n b_{-i}t^{-i}dt$.
  Grading by $\deg(a_i)=1$, $\deg(b_i)=0$ encodes a $\Gm$-action, and
  Proposition~\texttt{p:zn-coords} gives a $\Gm$-equivariant isomorphism
  $\Spec(A_n) \simeq \widetilde{\cZ}^{\leq n}$ with
  $\Spec(A_n)/\Gm \simeq \cZ^{\leq n}$. For $n=0$, $A_0 = k[a_0]$ is
  Noetherian; for $n=1$, $A_1 = k[b_{-1}, a_0, a_1, \ldots]/((b_{-1}-i)a_i)$;
  for general $n$ the ring $A_n$ is \emph{non-Noetherian}.

  \textit{Source: \cite{hilburn-raskin}, \S\texttt{ss:coords}.}
\end{definition}

\section{Local Abel-Jacobi morphisms}

\begin{proposition}[Log Contou-Carr\`ere pairing]
  \label{prop:cc-log}
  \lean{MR4583777TatesThesisInTheDeRhamSetting.TODO.ccLogPairing}
  \uses{def:positive-loop-space}
  % SOURCE: Hilburn-Raskin Proposition \ref{p:cc-log}, §\ref{s:aj}
  % SOURCE QUOTE: "Each of the morphisms: $\Gamma^{\IndCoh}(\Gr_{\bG_m}^{pos},\omega_{\Gr_{\bG_m}^{pos}}) \to \Gamma(\fL^+ \bA^1,\sO_{\fL^+ \bA^1})$ [and its $\Gr_{\bG_m}$, truncated, and $\Gr_{\bG_m}^{\leq n}$ variants] coming from our pairings and \S \ref{ss:pairing-map} is an isomorphism. Moreover, each isomorphism is a map of commutative algebras..."
  The bilinear ``log Contou-Carr\`ere'' pairings induce isomorphisms of
  commutative algebras
  \[
    \Gamma^{\IndCoh}(\GrGmpos, \omega_{\GrGmpos}) \xrightarrow{\sim}
    \Gamma(\LLAp, \cO_{\LLAp}),
    \quad
    \Gamma^{\IndCoh}(\GrGm, \omega_{\GrGm}) \xrightarrow{\sim}
    \Gamma(\LLGmp, \cO_{\LLGmp}),
  \]
  together with their truncated counterparts for $\GrGmposn$ and $\GrGmn$.
  The pairing $\pair{-}{-}:\GrGmpos \times \LLAp \to \BA^1$ restricts on
  $\GrGmpos \times \LLGmp$ to a pairing into $\Gm$ that agrees with the
  classical Contou-Carr\`ere pairing $\pair{-}{-}:\GrGm \times \LLGmp \to \Gm$
  of \cite{cc}.

  \textit{Source: \cite{hilburn-raskin}, Proposition \texttt{p:cc-log}.}
\end{proposition}

\section{Coherent sheaves on $\cY$}
\label{sec:indcoh}

\begin{definition}[Coherent sheaves on an eventually coconnective affine]
  \label{def:coh-S}
  \lean{MR4583777TatesThesisInTheDeRhamSetting.TODO.cohS}
  \uses{}
  % SOURCE: Hilburn-Raskin §\ref{s:indcoh} (Affine case)
  % SOURCE QUOTE: "We define the (non-cocomplete, DG) subcategory $\Coh(S) \subset \QCoh(S)$ to consist of objects $\sF \in \QCoh(S)^+$ such that for each $n$, the object $\tau^{\geq -n}(\sF) \in \QCoh(S)^{\geq -n}$ is compact."
  Let $S$ be an eventually coconnective (e.g.\ classical) affine scheme.
  The (non-cocomplete) DG subcategory
  \[
    \CohS(S) \;\subset\; \QCohS(S)
  \]
  consists of those $\sF \in \QCohS(S)^+$ such that for each $n$ the
  truncation $\tau^{\geq -n}(\sF) \in \QCohS(S)^{\geq -n}$ is compact;
  equivalently, $\sF$ is eventually coconnective and
  $\underline{\Hom}(\sF,-)$ commutes with filtered colimits in each
  $\QCohS(S)^{\geq -n}$. One has $\on{Perf}(S) \subset \CohS(S)$.

  \textit{Source: \cite{hilburn-raskin}, \S\texttt{s:indcoh}.}
\end{definition}

\begin{definition}[Ind-coherent sheaves $\IndCohStar(S)$]
  \label{def:indcoh-star-S}
  \lean{MR4583777TatesThesisInTheDeRhamSetting.TODO.indCohStarS}
  \uses{def:coh-S}
  % SOURCE: Hilburn-Raskin §\ref{s:indcoh} (Affine case)
  % SOURCE QUOTE: "We define $\IndCoh^*(S)$ as $\Ind(\Coh(S))$. Observe that we have a natural functor $\Psi:\IndCoh^*(S) \to \QCoh(S)$."
  For an eventually coconnective affine scheme $S$, define
  \[
    \IndCohStar(S) \;:=\; \Ind(\CohS(S)),
  \]
  the ind-completion of $\CohS(S)$. There is a canonical functor
  $\Psi : \IndCohStar(S) \to \QCohS(S)$ extending the inclusion
  $\CohS(S) \hookrightarrow \QCohS(S)$ by colimits.

  \textit{Source: \cite{hilburn-raskin}, \S4.2.1.}
\end{definition}

\begin{lemma}[$\Psi$ is an equivalence on bounded-below objects]
  \label{lem:psi-equiv}
  \lean{MR4583777TatesThesisInTheDeRhamSetting.TODO.psiEquiv}
  \uses{def:indcoh-star-S}
  % SOURCE: Hilburn-Raskin Lemma \ref{l:psi-affine} (citing \cite{methods} Lemma 6.4.1)
  % SOURCE QUOTE: "The above functor $\Psi$ is $t$-exact and induces an equivalence $\Psi:\IndCoh^*(S)^+ \isom \QCoh(S)^+$."
  The functor $\Psi : \IndCohStar(S) \to \QCohS(S)$ is $t$-exact and induces
  an equivalence
  \[
    \Psi : \IndCohStar(S)^+ \;\simeq\; \QCohS(S)^+
  \]
  on the bounded-below (eventually coconnective) full subcategories.

  \textit{Source: \cite{hilburn-raskin}, Lemma \texttt{l:psi-affine}, citing \cite{gaitsgory-rozenblyum-methods}, Lemma 6.4.1.}
\end{lemma}

\begin{proof}
  \uses{def:indcoh-star-S}
  This is Lemma~6.4.1 of Gaitsgory--Rozenblyum \cite{gaitsgory-rozenblyum-methods}.
  $t$-exactness of $\Psi$ follows because $\Psi$ is the ind-extension of
  the $t$-exact inclusion $\CohS(S) \hookrightarrow \QCohS(S)$.
  For essential surjectivity on bounded-below objects: given
  $M \in \QCohS(S)^+$, write $M = \varinjlim_r \tau^{\geq -r} M$.
  Each $\tau^{\geq -r} M$ is bounded below, and by the eventually
  coconnective hypothesis on $S$, every bounded-below quasi-coherent
  sheaf is a filtered colimit of objects in $\CohS(S)$ (using
  truncation $\tau^{\geq -r}$ and the fact that $\CohS(S)$ generates
  $\QCohS(S)^+$). Hence $M \in \Psi(\IndCohStar(S)^+)$, giving
  essential surjectivity. Full faithfulness on bounded-below objects
  follows from the universal property of $\Ind(\CohS(S))$.
\end{proof}

\begin{proposition}[t-structure on $\IndCohStar(S)$]
  \label{prop:indcoh-t-structure}
  \lean{MR4583777TatesThesisInTheDeRhamSetting.TODO.indCohTStructure}
  \uses{def:indcoh-star-S, lem:psi-equiv}
  % SOURCE: Hilburn-Raskin §\ref{s:indcoh} (Affine case)
  % SOURCE QUOTE: "Define $\IndCoh^*(S)^{\leq 0}$ as the subcategory generated under colimits by coherent objects that are connective (with respect to the $t$-structure on $\QCoh(S)$). There is a corresponding $t$-structure on $\IndCoh^*(S)$ with this as its connective subcategory."
  Defining $\IndCohStar(S)^{\leq 0}$ as the subcategory generated under
  colimits by connective coherent objects yields a $t$-structure on
  $\IndCohStar(S)$ for which $\Psi$ is $t$-exact (Lemma~\ref{lem:psi-equiv}).
  Its heart is the abelian category of coherent $H^0(\mathcal{O}(S))$-modules.

  \textit{Source: \cite{hilburn-raskin}, \S\texttt{s:indcoh}.}
\end{proposition}

\begin{proof}
  \uses{def:indcoh-star-S, lem:psi-equiv}
  Existence: define $\IndCohStar(S)^{\leq 0}$ to be the preimage of
  $\QCohS(S)^{\leq 0}$ under $\Psi$. Since $\Psi$ is $t$-exact
  (Lemma~\ref{lem:psi-equiv}), it intertwines the truncation functors
  $\tau^{\leq 0}$ on $\IndCohStar(S)$ and $\QCohS(S)$, and the axioms
  of a $t$-structure (existence of left and right truncations, the
  half-plane conditions) follow from those on $\QCohS(S)$ by transport
  along $\Psi$. Uniqueness is clear since $\Psi$ is $t$-exact and
  essentially surjective on bounded-below objects (Lemma~\ref{lem:psi-equiv}).
  The identification of the heart with coherent $H^0(\mathcal{O}(S))$-modules
  follows from the heart of the standard $t$-structure on $\QCohS(S)$.
\end{proof}

\begin{definition}[Semi-coherent algebra]
  \label{def:semi-coherent}
  \lean{MR4583777TatesThesisInTheDeRhamSetting.TODO.semiCoherent}
  \uses{def:coh-S}
  % SOURCE: Hilburn-Raskin §\ref{s:indcoh} (Semi-coherence), Definition
  % SOURCE QUOTE: "A connective and eventually coconnective commutative algebra $A$ is \emph{semi-coherent} if every $M \in A\mod^+$ can be written as a filtered colimit of objects $M_i$ such that: $M_i \in \Coh(A)$ for every $i$. There exists an integer $r$ independent of $i$ such that every $M_i$ lies in $A\mod^{\geq -r}$."
  A connective and eventually coconnective commutative algebra $A$ is
  \emph{semi-coherent} if every $M \in A\mbox{-}\on{mod}^+$ (bounded-below
  $A$-modules) can be written as a filtered colimit of objects $M_i$ such that:
  \begin{itemize}
    \item $M_i \in \CohS(A)$ for every $i$, and
    \item there exists an integer $r$ independent of $i$ such that every
          $M_i \in A\mbox{-}\on{mod}^{\geq -r}$.
  \end{itemize}
  ($S = \Spec(A)$ is called semi-coherent if $A$ is.)

  \textit{Source: \cite{hilburn-raskin}, \S\texttt{s:indcoh} (Semi-coherence).}
\end{definition}

\begin{proposition}[Semi-coherence of $A_n$]
  \label{prop:An-semi-coherent}
  \lean{MR4583777TatesThesisInTheDeRhamSetting.TODO.AnSemiCoherent}
  \uses{def:An-coords, def:semi-coherent}
  % SOURCE: Hilburn-Raskin Proposition \ref{p:semi-coh}
  % SOURCE QUOTE: "The (classical) commutative algebra $A_n$ defined in \S \ref{ss:coords} is semi-coherent." / "In other words, $\widetilde{\sZ}^{\leq n}$ is semi-coherent."
  The classical commutative algebra $A_n$ defined in
  Definition~\ref{def:An-coords} is semi-coherent; equivalently,
  $\widetilde{\cZ}^{\leq n}$ is semi-coherent. (Note: $A_n$ is
  \emph{not} Noetherian for $n \geq 1$, so this is a genuine theorem,
  not a formal consequence of Noetherianity.)

  \textit{Source: \cite{hilburn-raskin}, Proposition \texttt{p:semi-coh}.}
\end{proposition}

% SOURCE QUOTE PROOF: (read from references/MR4583777-de-rham-tate.tex, \S \ref{ss:z-semi-coh-pf})
% "Our argument is by induction on $n$. We proceed in steps."
% Step 1 ($n=0$): "$A_0 = k[a_0]$ is Noetherian."
% Step 2 ($n=1$): "$A_1$ is actually coherent." (via the Noetherian subalgebras
%   $A_{1,\leq m} \subset A_1$: any f.p. $M$ descends to some $A_{1,\leq m}$.)
% Step 3: define "$r$-good" = filtered colimit of coherent objects in degrees $\geq -r$;
%   $r$-good objects are closed under filtered colimits and extensions.
% Step 4: for $n>1$ assume the result for $n-1$; reduce to showing $M[b_{-n}^{-1}]$
%   and $\Coker(M \to M[b_{-n}^{-1}])$ are good.
% Step 5 (generic case): $A_n[b_{-n}^{-1}] = k[b_{-1},\ldots,b_{-n},b_{-n}^{-1}]$ is
%   regular Noetherian, so $M[b_{-n}^{-1}]$ has bounded Tor-amplitude, hence (Lazard) is good.
% Step 6 (torsion case): for $\widetilde{M}$ with $b_{-n}$ locally nilpotent, filter by
%   kernels of $b_{-n}^i$; quotients are $A_{n-1}$-modules, $r$-good by induction (uniform $r$),
%   so $\widetilde{M}$ is $r$-good.
% Step 7: the $n=1$ case can be redone in this framework by localizing at
%   $\{b_{-1}, b_{-1}-1, \ldots\}$.
\begin{proof}
  \uses{def:An-coords, def:semi-coherent}
  The proof (\cite{hilburn-raskin}, \S\texttt{ss:z-semi-coh-pf}) is by
  induction on $n$ and proceeds in seven steps.
  \emph{Step 1} ($n=0$): $A_0 = k[a_0]$ is Noetherian, hence semi-coherent.
  \emph{Step 2} ($n=1$): one shows directly that $A_1$ is in fact coherent,
  using the ascending chain of \emph{Noetherian} subalgebras
  $A_{1,\leq m} = k[b_{-1}, a_0, \ldots, a_m] \subset A_1$: any finitely
  presented $M \in A_1\text{-}\on{mod}^\heartsuit$ is base-changed from some
  finitely generated module over $A_{1,\leq m}$ for $m \gg 0$.
  \emph{Step 3} (d\'evissage): call $M \in A_n\text{-}\on{mod}$ \emph{$r$-good}
  if it is a filtered colimit of coherent objects concentrated in degrees
  $\geq -r$ (equivalently $M \in \Ind(\CohS(A_n)\cap A_n\text{-}\on{mod}^{\geq -r})$),
  and \emph{good} if $r$-good for some $r$; semi-coherence amounts to every
  $M \in A_n\text{-}\on{mod}^{\heartsuit}$ being good. The $r$-good objects
  (fixed $r$) are closed under filtered colimits and extensions.
  \emph{Step 4} (reduction): for $n > 1$, assuming the result for $n-1$, it
  suffices to show that $M[b_{-n}^{-1}] := \varinjlim_{b_{-n}} M$ and the
  homotopy cokernel $\Coker(M \to M[b_{-n}^{-1}])$ are good.
  \emph{Step 5} (generic case): $A_n[b_{-n}^{-1}] = k[b_{-1},\ldots,b_{-n},b_{-n}^{-1}]$
  is regular Noetherian, so $M[b_{-n}^{-1}]$ has bounded $\on{Tor}$-amplitude
  over $A_n$; by Lazard's theorem a flat $A_n$-module is good, so
  $M[b_{-n}^{-1}]$ is good.
  \emph{Step 6} (torsion case): if $b_{-n}$ acts locally nilpotently on
  $\widetilde{M}$, filter by the kernels $\widetilde{M}_i$ of $b_{-n}^i$; the
  successive quotients $\widetilde{M}_i/\widetilde{M}_{i-1}$ are
  $A_{n-1}$-modules, hence $r$-good (with $r$ independent of $i$) by the
  inductive hypothesis and the almost-finite-presentation of
  $A_n \to A_{n-1}$, so $\widetilde{M} = \varinjlim_i \widetilde{M}_i$ is $r$-good.
  \emph{Step 7}: the $n=1$ case fits the same framework by localizing $A_1$ at
  $\{b_{-1}, b_{-1}-1, b_{-1}-2, \ldots\}$ to reach a regular ring, with the
  goodness bound for torsion modules uniform in the element.
\end{proof}

\begin{definition}[$\IndCohStar(\cY)$ via Grassmannian weak-invariants]
  \label{def:indcoh-star-Y}
  \lean{MR4583777TatesThesisInTheDeRhamSetting.TODO.indCohStarY}
  \uses{def:indcoh-star-S, def:Y, def:Z-leq-n}
  % SOURCE: Hilburn-Raskin §4.7 (Ind-coherent sheaves on $\sY$)
  % SOURCE QUOTE: "We now define: $\IndCoh^*(\sY) \coloneqq \IndCoh^*(\sY_{\log})^{\Gr_{\bG_m},w}$. By definition, the right hand side is: $\TwoHom_{\IndCoh(\Gr_{\bG_m})\mod}(\Vect, \IndCoh^*(\sY_{\log}))$."
  Recall $\cY = \cY_{\log}/\GrGm$, with $\IndCoh(\GrGm)$ acting on
  $\IndCohStar(\cY_{\log})$. Define $\IndCohStar(\cY)$ as the \emph{weak}
  ($\GrGm$-)invariants
  \[
    \IndCohStar(\cY) \;:=\; \IndCohStar(\cY_{\log})^{\GrGm,\, w}
    \;:=\; \on{Hom}_{\IndCoh(\GrGm)\text{-}\on{mod}}
           \big(\Vect,\, \IndCohStar(\cY_{\log})\big).
  \]
  There is a conservative forgetful functor
  $\Oblv : \IndCohStar(\cY) \to \IndCohStar(\cY_{\log})$, with left adjoint
  $\Av_!^{\GrGm,\, w}$. In the truncated setting one similarly defines
  $\IndCohStar(\cY^{\leq n}) := \IndCohStar(\cY_{\log}^{\leq n})^{\GrGmn,\, w}$.

  \textit{Source: \cite{hilburn-raskin}, \S4.7.}
\end{definition}

\begin{lemma}[Full faithfulness of $\lambda_{n,*}^{\IndCoh}$]
  \label{lem:lambda-n-ff}
  \lean{MR4583777TatesThesisInTheDeRhamSetting.TODO.lambdaNFullyFaithful}
  \uses{def:indcoh-star-Y, def:Z-leq-n}
  % SOURCE: Hilburn-Raskin Lemma \ref{l:lambda-ff}
  % SOURCE QUOTE: "For $n>0$, the functor $\lambda_{n,*}^{\IndCoh}$ is fully faithful."
  % (where $\lambda_{n,*}^{\IndCoh}:\IndCoh^*(\sY^{\leq n}) \to \IndCoh^*(\sY)$ is
  % $t$-exact and preserves compact objects.)
  For $n > 0$, the $\IndCoh$-\emph{pushforward} functor
  \[
    \lambda_{n,*}^{\IndCoh} : \IndCohStar(\cY^{\leq n})
    \;\longrightarrow\; \IndCohStar(\cY)
  \]
  (which is $t$-exact and preserves compact objects) is fully faithful.

  \textit{Source: \cite{hilburn-raskin}, Lemma \texttt{l:lambda-ff}.}
\end{lemma}

% SOURCE QUOTE PROOF: (read from references/MR4583777-de-rham-tate.tex)
% "Roughly speaking, this is true because the map $\sY^{\leq n} \to \sY$ is
% formally \'etale for $n>0$, which follows from the corresponding
% fact for $\LocSys_{\bG_m}^{\leq n} \to \LocSys_{\bG_m}$ (which
% follows from Proposition \ref{p:locsys-log}). Unwinding the constructions
% to convert this argument into a proof is straightforward."
\begin{proof}
  \uses{def:indcoh-star-Y, def:Z-leq-n, prop:locsys-log-decomp}
  Roughly speaking, this holds because the map $\cY^{\leq n} \to \cY$ is
  formally \'etale for $n > 0$, which follows from the corresponding fact
  for the truncated comparison $\LocSysGm^{\leq n} \to \LocSysGm$
  (Proposition~\ref{prop:locsys-log-decomp}, i.e.\ Proposition~\texttt{p:locsys-log}
  in \cite{hilburn-raskin}). Unwinding the constructions to convert this
  argument into a proof is straightforward.
\end{proof}

\section{Spectral realization of Weyl algebras}
\label{sec:dmod}

\begin{definition}[Weyl algebra $W_n$]
  \label{def:weyl-algebra-Wn}
  \lean{MR4583777TatesThesisInTheDeRhamSetting.TODO.weylAlgebraWn}
  \uses{def:truncated-positive-loop-space}
  % SOURCE: Hilburn-Raskin §\ref{s:weyl}
  % SOURCE QUOTE: "let $W_n$ denote the algebra of global differential operators on the scheme $\fL_n^+ \bA^1 = \Spec\Sym(t^{-n}k[[t]]dt/k[[t]]dt)$. Let $W_n^{op}$ denote $W_n$ with the reversed multiplication."
  The \emph{Weyl algebra} $W_n$ is the algebra of global differential operators
  on the scheme $\LLAn{n} \cong \BA^n$:
  \[
    W_n \;:=\; k\langle x_0,\ldots,x_{n-1},\, \partial_0,\ldots,\partial_{n-1}
    \rangle \;\big/\; ([\partial_i, x_j] = \delta_{ij},\; [x_i,x_j] = 0,\;
    [\partial_i,\partial_j] = 0).
  \]
  As a $k$-algebra, $W_n \cong D(\BA^n)$ where $D$ denotes the ring of
  algebraic differential operators; $W_n^{op}$ denotes $W_n$ with the
  reversed multiplication.

  \textit{Source: \cite{hilburn-raskin}, \S\texttt{s:weyl}.}
\end{definition}

\begin{theorem}[Spectral realization of the Weyl algebra]
  \label{thm:spectral-realization}
  \lean{MR4583777TatesThesisInTheDeRhamSetting.TODO.spectralRealization}
  \uses{def:weyl-algebra-Wn, def:indcoh-star-Y, def:Z-leq-n, lem:lambda-n-ff}
  % SOURCE: Hilburn-Raskin §\ref{s:weyl}, eqs \eqref{eq:action} and Prop \ref{p:ff}
  % SOURCE QUOTE: "Let $\sF_n \coloneqq \Av_!^{\Gr_{\bG_m}^{\leq n},w}(i_{n,*}(\sO_{\sZ^{\leq n}}))$, which is a compact object in $\IndCoh^*(\sY^{\leq n})$. ... We will construct a canonical homomorphism: $W_n^{op} \to \ul{\End}_{\IndCoh^*(\sY^{\leq n})}(\sF_n) \in \Alg = \Alg(\Vect)$." / Prop \ref{p:ff}: "we will show that this map is actually an isomorphism."
  For each $n \geq 0$, let
  $\cF_n := \Av_!^{\GrGmn,\, w}(i_{n,*}(\cO_{\cZ^{\leq n}}))$, a compact object
  of $\IndCohStar(\cY^{\leq n})$ (where $i_n : \cZ^{\leq n}\hookrightarrow
  \cY_{\log}^{\leq n}$ is the canonical embedding). There is a canonical
  algebra homomorphism
  \[
    W_n^{op} \;\longrightarrow\;
    \underline{\on{End}}_{\IndCohStar(\cY^{\leq n})}(\cF_n),
  \]
  and (Proposition~\texttt{p:ff}) it is an isomorphism.

  \textit{Source: \cite{hilburn-raskin}, \S\texttt{s:weyl}, equation \texttt{eq:action} and Proposition \texttt{p:ff}.}
\end{theorem}

% SOURCE QUOTE PROOF: (read from references/MR4583777-de-rham-tate.tex)
% "Let $\sF_n \coloneqq \Av_!^{\Gr_{\bG_m}^{\leq n},w}
% (i_{n,*}(\sO_{\sZ^{\leq n}}))$, which is a compact object in
% $\IndCoh^*(\sY^{\leq n})$.
% In this section, we construct an action of a Weyl algebra in $2n$
% generators on $\sF_n$ (considered as an object of $\IndCoh^*(\sY^{\leq n})$).
% We will construct a canonical homomorphism:
% $W_n^{op} \to \ul{\End}_{\IndCoh^*(\sY^{\leq n})}(\sF_n) \in \Alg = \Alg(\Vect)$."
\begin{proof}
  \uses{def:weyl-algebra-Wn, def:indcoh-star-Y, def:Z-leq-n, lem:lambda-n-ff,
        prop:Z-leq-n-algebraic, prop:indcoh-t-structure}
  Set $\cF_n := \Av_!^{\GrGmn,\, w}(i_{n,*}(\cO_{\cZ^{\leq n}}))$, a compact
  object of $\IndCohStar(\cY^{\leq n})$ lying in the heart of the $t$-structure
  of Proposition~\ref{prop:indcoh-t-structure}; here
  $i_n : \cZ^{\leq n} \hookrightarrow \cY_{\log}^{\leq n}$ is the canonical
  embedding. Since $\cF_n$ lies in the heart, the target
  $\underline{\on{End}}_{\IndCohStar(\cY^{\leq n})}(\cF_n)$ is concentrated in
  degrees $\geq 0$, so it suffices to build the homomorphism into
  $H^0\underline{\on{End}}_{\IndCohStar(\cY^{\leq n})}(\cF_n)$ by generators
  and relations. One constructs:
  \begin{itemize}
    \item maps $\omega \mapsto \varphi_\omega$ from
      $t^{-n}k[[t]]dt/k[[t]]dt = (k[[t]]/t^n)^\vee$, and
    \item maps $f \mapsto \xi_f$ from $k[[t]]/t^n$,
  \end{itemize}
  and verifies that the $\varphi_\omega$ mutually commute, the $\xi_f$
  mutually commute, and the canonical commutation relation
  $[\xi_f, \varphi_\omega] = -\Res(f\omega)\cdot\on{id}$ holds (the sign
  reflecting the use of $W_n^{op}$). That this homomorphism is an
  \emph{isomorphism} is Proposition~\texttt{p:ff} in \cite{hilburn-raskin},
  whose proof is by induction on $n$ using the Contou-Carr\`ere pairing.
\end{proof}

\section{Compatibility with class field theory}
\label{sec:lcft}

This section records the local geometric class field theory (Contou-Carr\`ere)
content that is part of the statement of the main theorem: the equivalence
$\Delta$ is asserted to be one of $D^*(\LLGm)$-module categories.

\begin{theorem}[Truncated local class field theory]
  \label{thm:trun-lcft}
  \lean{MR4583777TatesThesisInTheDeRhamSetting.TODO.trunLCFT}
  \uses{prop:cc-log, def:locsys-gm}
  % SOURCE: Hilburn-Raskin Theorem \ref{t:trun-lcft}, §\ref{s:cft}
  % SOURCE QUOTE: "The functor [\eqref{eq:trun-lcft}: $D(\fL\bG_m/\sK_n) \to \QCoh(\LocSys_{\bG_m}^{\leq n})$] is an equivalence. In particular, there is a canonical fully faithful symmetric monoidal functor $D(\fL_n^+\bG_m) \into \QCoh(\LocSys_{\bG_m}^{\leq n})$."
  Let $\sK_n := \on{Ker}(\LLGmp \to \fL^+_n\Gm)$ be the $n$th congruence
  subgroup. The Fourier--Mukai functor
  $D(\LLGm/\sK_n) \to \QCohS(\LocSysGm^{\leq n})$ arising from the Cartier
  duality pairing is a symmetric monoidal equivalence. In particular, there
  is a canonical fully faithful symmetric monoidal functor
  $D(\fL^+_n\Gm) \hookrightarrow \QCohS(\LocSysGm^{\leq n})$.

  \textit{Source: \cite{hilburn-raskin}, Theorem \texttt{t:trun-lcft}.}
\end{theorem}

\begin{theorem}[Full local class field theory of Beilinson--Drinfeld]
  \label{thm:full-lcft}
  \lean{MR4583777TatesThesisInTheDeRhamSetting.TODO.fullLCFT}
  \uses{thm:trun-lcft, def:locsys-gm}
  % SOURCE: Hilburn-Raskin Theorem \ref{t:full-lcft}, §\ref{s:cft}
  % SOURCE QUOTE: "There is a canonical symmetric monoidal equivalence: $D^*(\fL\bG_m) \isom \QCoh(\LocSys_{\bG_m})$."
  Passing to the limit over $n$ yields the Beilinson--Drinfeld equivalence: a
  canonical symmetric monoidal equivalence
  \[
    D^*(\LLGm) \;\xrightarrow{\;\sim\;}\; \QCohS(\LocSysGm).
  \]

  \textit{Source: \cite{hilburn-raskin}, Theorem \texttt{t:full-lcft}.}
\end{theorem}

\begin{proposition}[$\overline{\Delta}_n$ is $D(\fL^+_n\Gm)$-equivariant]
  \label{prop:lcft-compat}
  \lean{MR4583777TatesThesisInTheDeRhamSetting.TODO.deltaLCFTEquivariant}
  \uses{thm:trun-lcft, def:indcoh-star-Y, def:dmod-loop-space}
  % SOURCE: Hilburn-Raskin §\ref{s:cft}, "Formulation of the equivariance property"
  % SOURCE QUOTE: "Now observe that $\ol{\Delta}_n$ maps a category acted on by $D(\fL_n^+\bG_m)$ to one acted on by $\QCoh(\LocSys_{\bG_m}^{\leq n})$ (via the structure map $\sY^{\leq n} \to \LocSys_{\bG_m}^{\leq n}$). Using Theorem \ref{t:trun-lcft}, we can regard both sides as acted on by $D(\fL_n^+\bG_m)$. ... we show that $\ol{\Delta}_n$ is canonically a morphism of $D(\fL_n^+\bG_m)$-module categories."
  The source $D(\LLAn{n}) \simeq W_n\text{-}\on{mod}$ of $\overline{\Delta}_n$
  is acted on by $D(\fL^+_n\Gm)$, while the target $\IndCohStar(\cY^{\leq n})$
  is acted on by $\QCohS(\LocSysGm^{\leq n})$ via the structure map
  $\cY^{\leq n} \to \LocSysGm^{\leq n}$. Identifying the two actions through
  Theorem~\ref{thm:trun-lcft}, the functor $\overline{\Delta}_n$ is
  canonically a morphism of $D(\fL^+_n\Gm)$-module categories. This
  equivariance is the form in which local geometric class field theory enters
  the inductive proof of full faithfulness.

  \textit{Source: \cite{hilburn-raskin}, \S\texttt{s:cft}.}
\end{proposition}

\section{Fully faithfulness}

\begin{theorem}[Main functor is fully faithful]
  \label{thm:fully-faithful}
  \lean{MR4583777TatesThesisInTheDeRhamSetting.TODO.mainFunctorFullyFaithful}
  \uses{thm:spectral-realization, def:dmod-loop-space, def:indcoh-star-Y,
        prop:indcoh-t-structure, lem:lambda-n-ff, prop:lcft-compat}
  % SOURCE: Hilburn-Raskin Proposition \ref{p:ff}
  % SOURCE QUOTE: "For every $n \geq 0$, the functor $\Delta_n$ is fully faithful."
  For each $n \geq 0$ the functor
  $\Delta_n : D(\LLAn{n}) \simeq W_n\text{-}\on{mod} \to \IndCohStar(\cY)$
  (the composite $\overline{\Delta}_n$ followed by $\lambda_{n,*}^{\IndCoh}$)
  is fully faithful, and hence so is the colimit
  \[
    \Delta \;=\; \varinjlim_n \Delta_n \;:\; D^!(\LLA) \;\longrightarrow\;
    \IndCohStar(\cY).
  \]

  \textit{Source: \cite{hilburn-raskin}, Proposition \texttt{p:ff} (with Lemma \texttt{l:colim-ff}).}
\end{theorem}

% SOURCE QUOTE PROOF: (read from references/MR4583777-de-rham-tate.tex)
% "Our proof is essentially by induction. We settle the $n = 0$ case
% by explicit analysis. We then use local class field theory
% (or the Contou-Carr\`ere pairing) to perform the inductive step."
\begin{proof}
  \uses{thm:spectral-realization, prop:indcoh-t-structure, lem:lambda-n-ff,
        def:dmod-loop-space, prop:lcft-compat}
  The argument is by induction on $n$. The functor $\Delta_n$ factors as
  $D(\LLAn{n}) \xrightarrow{\overline{\Delta}_n} \IndCohStar(\cY^{\leq n})
  \xrightarrow{\lambda_{n,*}^{\IndCoh}} \IndCohStar(\cY)$, where
  $\overline{\Delta}_n$ sends $D_{\LLAn{n}}$ to $\cF_n$ compatibly with the
  spectral realization $W_n^{op} \xrightarrow{\sim}
  \underline{\on{End}}(\cF_n)$ (Theorem~\ref{thm:spectral-realization}).
  By Lemma~\ref{lem:lambda-n-ff}, $\lambda_{n,*}^{\IndCoh}$ is fully faithful
  for $n > 0$, so it suffices to prove $\overline{\Delta}_n$ fully faithful
  (Proposition~\texttt{p:ff}). The base case $n = 0$ is by explicit analysis;
  the inductive step uses the local geometric class field theory
  (Contou-Carr\`ere) equivariance of $\overline{\Delta}_n$
  (Proposition~\ref{prop:lcft-compat}) to compare the $n+1$ level with the
  $n$ level. Full faithfulness of $\Delta = \varinjlim_n \Delta_n$ then
  follows from Lemma~\texttt{l:colim-ff} in \cite{hilburn-raskin}: a filtered
  colimit of fully faithful, compact-object-preserving functors between
  compactly generated categories remains fully faithful.
\end{proof}

\section{Proof of the main theorem}
\label{sec:main-results}

\begin{definition}[D-module category on the loop space]
  \label{def:dmod-loop-space}
  \lean{MR4583777TatesThesisInTheDeRhamSetting.TODO.dmodLoopSpace}
  \uses{def:loop-space, def:truncated-positive-loop-space, def:weyl-algebra-Wn}
  % SOURCE: Hilburn-Raskin §\ref{s:main} (citing \cite{dmod}), and proof of Thm \ref{t:main-precise}
  % SOURCE QUOTE: "Recall from \cite{dmod} that by definition we have: $D^!(\fL^+\bA^1) \coloneqq \colim_{n \geq 0} \, D(\fL_n^+\bA^1)$ where the structural functors are the standard $!$-pullback functors of $D$-module theory."
  % SOURCE QUOTE (full loop space): "$D^*(\fL\bG_m) \underset{D^*(\fL^{pos}\bG_m)}{\otimes} D^!(\fL^+\bA^1) \to D^!(\fL\bA^1)$ ... is an equivalence."
  The \emph{D-module category on the positive loop space} is the colimit
  \[
    D^!(\LLAp) \;:=\; \varinjlim_{n \geq 0}\; D(\LLAn{n})
    \;\simeq\; \varinjlim_n\; W_n\mbox{-}\on{mod},
  \]
  taken along the standard $!$-pullback transition functors
  $D(\LLAn{n}) \to D(\LLAn{n+1})$ of $D$-module theory. The category for the
  \emph{full} loop space is obtained by class-field-theory induction,
  \[
    D^!(\LLA) \;\simeq\; D^*(\LLGm)
    \underset{D^*(\fL^{pos}\Gm)}{\otimes} D^!(\LLAp),
  \]
  using a coordinate $t$ to split $\LLGm \simeq \LLGmp \times \BZ$.

  \textit{Source: \cite{hilburn-raskin}, \S\texttt{s:main} (citing \cite{dmod}).}
\end{definition}

\begin{theorem}[Main equivalence]
  \label{thm:main}
  \lean{MR4583777TatesThesisInTheDeRhamSetting.TODO.mainEquivalence}
  \uses{thm:fully-faithful, thm:Y-classical, prop:Z-leq-n-algebraic,
        def:indcoh-star-Y, def:dmod-loop-space,
        prop:An-semi-coherent, prop:locsys-decomp,
        thm:full-lcft, prop:lcft-compat}
  % SOURCE: Hilburn-Raskin Theorem \ref{t:main-precise} (= Theorem \ref{t:main} of the intro)
  % SOURCE QUOTE: "There is a unique equivalence: $\Delta:D^!(\fL\bA^1) \isom \IndCoh^*(\sY)$ of $D^*(\fL\bG_m)$-module categories such that the induced morphism: $D^!(\fL^+\bA^1) \subset D^!(\fL\bA^1) \to \IndCoh^*(\sY)$ of $D^*(\fL^{pos}\bG_m)$-module categories coincides with the functor $\Delta_{\infty}$ from above."
  % INTRO QUOTE: "There is a canonical equivalence of DG categories: $\Delta:D^!(\fL \bA^1) \simeq \IndCoh^*(\sY)$. Moreover, this equivalence is compatible with the local geometric class field theory of Beilinson-Drinfeld."
  There is a unique equivalence of DG categories
  \[
    \Delta \;:\; D^!(\LLA) \;\xrightarrow{\;\sim\;}\; \IndCohStar(\cY)
  \]
  that is a morphism of $D^*(\LLGm)$-module categories, where $D^*(\LLGm)$
  acts on the left via class field theory and on the right (through
  Theorem~\ref{thm:full-lcft}, $D^*(\LLGm)\simeq\QCohS(\LocSysGm)$) via the
  structure map $\cY \to \LocSysGm$, and which restricts on
  $D^!(\LLAp)\subset D^!(\LLA)$ to the functor $\Delta_\infty$. This is the
  de Rham analogue of Tate's thesis: a canonical equivalence compatible with
  the local geometric class field theory of Beilinson--Drinfeld.

  \textit{Source: \cite{hilburn-raskin}, Theorem \texttt{t:main-precise} (= Theorem \texttt{t:main}).}
\end{theorem}

% SOURCE QUOTE PROOF: (read from references/MR4583777-de-rham-tate.tex)
% "First, we note that $\Delta_{\infty}$ is fully faithful by
% Proposition \ref{p:ff} and Lemma \ref{l:colim-ff} below. Next,
% we deduce that $\Delta$ itself is fully faithful by another
% application of Lemma \ref{l:colim-ff}.
% Therefore, it remains to show that $\Delta$ is essentially surjective.
% Clearly its essential image contains each object:
% $\pi_{n,*}^{\IndCoh}(\sO_{\sZ^{\leq n}}(m))$
% for all $n \geq 0$, $m \in \bZ$. Therefore, it suffices
% to show that $\IndCoh^*(\sY)$ is generated under colimits by these objects.
% But this follows immediately from the definitions and Corollary \ref{c:z-gen}."
\begin{proof}
  \uses{thm:fully-faithful, thm:Y-classical, prop:Z-leq-n-algebraic,
        prop:An-semi-coherent, thm:spectral-realization, def:indcoh-star-Y,
        thm:full-lcft, prop:lcft-compat}
  First, the functor $D^*(\LLGm)\otimes_{D^*(\fL^{pos}\Gm)} D^!(\LLAp)
  \to D^!(\LLA)$ is an equivalence (choosing a coordinate $t$ splits
  $\LLGm \simeq \LLGmp\times\BZ$), so there is a unique $D^*(\LLGm)$-module
  functor $\Delta$ restricting to $\Delta_\infty$. By
  Theorem~\ref{thm:fully-faithful} (Proposition~\texttt{p:ff} and
  Lemma~\texttt{l:colim-ff}) $\Delta_\infty$, hence $\Delta$, is fully
  faithful. For essential surjectivity, the essential image of $\Delta$
  contains every object $\pi_{n,*}^{\IndCoh}(\cO_{\cZ^{\leq n}}(m))$
  for $n \geq 0$, $m \in \BZ$, so it suffices that $\IndCohStar(\cY)$ be
  generated under colimits by these objects. This follows from the
  definitions together with Corollary~\texttt{c:z-gen} in \cite{hilburn-raskin},
  which rests on the classicality of $\cY$ (Theorem~\ref{thm:Y-classical}),
  the algebraicity of the $\cZ^{\leq n}$ (Proposition~\ref{prop:Z-leq-n-algebraic}),
  and the semi-coherence of $A_n$ (Proposition~\ref{prop:An-semi-coherent}).
  Compatibility with the $D^*(\LLGm)$-action is built into the construction
  of $\Delta$ via the class field theory equivalence
  (Theorem~\ref{thm:full-lcft}) and the equivariance of
  Proposition~\ref{prop:lcft-compat}.
\end{proof}

